%%%%%%%%%%%%%%%%%%%%%%%%%%%%%%%%%%%%%%%%%%%%%%%%%%%%%%%%%%%%%%%%%%%%%%
% Lớp 12 | Phương trình mặt phẳng, đường thẳng, mặt cầu | Phương trình mặt phẳng
% Dạng 2H5?1-2: Xác định véc-tơ pháp tuyến, cặp véc-tơ chỉ phương
% Số câu: 47  (sinh tự động bởi tach_cau.py; từ nay thầy sửa trực tiếp file này)
%%%%%%%%%%%%%%%%%%%%%%%%%%%%%%%%%%%%%%%%%%%%%%%%%%%%%%%%%%%%%%%%%%%%%%

%%%%%%%%%%%%%%%% TRẮC NGHIỆM 4 PHƯƠNG ÁN %%%%%%%%%%%%%%%%

% nguồn: Phan1_Lop_12_Hinh_hoc_Chuong_5
\begin{ex}%[2H5N1-2]%[ma:8b866063]
	Trong không gian $Oxyz$, vectơ nào sau đây~là vectơ pháp tuyến của mặt phẳng\xuongdong $(P)\colon x+3 y-4  z+5=0$?
	\choice 
		{$\vec{n}_{1}=(3; 4;  5)$}
		{\True  $\vec{n}_{2}=(1; 3;-4)$}
		{$\vec{n}_{3}=(1;  3; 4)$}
		 {$\vec{n}_{4}=(3; -4; 5)$}
	\loigiai{ 
		Dễ thấy rằng  mặt phẳng $(P)\colon x+3 y-4 z+5=0$ có một vectơ pháp tuyến~là $(1; 3; -4)$.
	} 
\end{ex}

% nguồn: Phan1_Lop_12_Hinh_hoc_Chuong_5
\begin{ex}%[2H5N1-2]%[ma:3d9526b1]
	Trong không gian $Oxyz$, vectơ nào sau đây là vectơ  pháp tuyến của mặt phẳng $(P)\colon x+3y-4z+5=0$?
	\choice 
		{\True $\overrightarrow{n}_2=(1; 3;-4)$} 
		{$\overrightarrow{n}_4=(3;-4;  5)$}
		 {$\overrightarrow{n}_3=(1; 3; 4)$}
		{$\overrightarrow{n}_1=(3;  4; 5)$}
	\loigiai{Vectơ pháp  tuyến của mặt phẳng $(P)\colon x+3y-4z+5=0$ là $\overrightarrow{n}_2=(1; 3;-4)$.}
\end{ex}

% nguồn: Phan1_Lop_12_Hinh_hoc_Chuong_5
\begin{ex}%[2H5N1-2]%[ma:b633c471]
	 Trong không gian $Oxyz$, mặt phẳng $(P)\colon x+y-2z+4=0$ có vectơ pháp tuyến là
	 \choice
		 {\True $(1;1;-2)$}
		{$(3;2;-2)$} 
		{$(1;1;2)$} 
		{$(2;1;-2)$} 
	\loigiai{ 
		Mặt phẳng $(P)\colon  x+y-2z+4=0$ có một vectơ pháp tuyến là $\overrightarrow{n}=(1;1;-2)$.
	} 
\end{ex}

% nguồn: Phan1_Lop_12_Hinh_hoc_Chuong_5
\begin{ex}%[2H5N1-2]%[ma:5d0fad61]
	Trong không gian $(Oxyz)$, cho mặt phẳng $(P)\colon x+2y-3z+6=0$. Véc-tơ nào sau đây là véc-tơ  pháp tuyến của mặt phẳng $(P)$?
	 \choice
		 {$\overrightarrow{n}=(-1; 2;-3)$}
		{$\overrightarrow{n}=(1;-2;  3)$}
		 {$\overrightarrow{n}=(-1;-2;-3)$}
		{\True $\overrightarrow{n}=(1; 2;-3)$} 
	 \loigiai{
		Ta có $\overrightarrow{n}=(1;2;-3)$ là véc-tơ pháp tuyến của  mặt phẳng $ (P) $.
	 }
\end{ex}

% nguồn: Phan1_Lop_12_Hinh_hoc_Chuong_5
\begin{ex}%[2H5N1-2]%[ma:8b933094]
	Trong không gian $Oxyz$,  cho mặt phẳng $(P)\colon x+y-z+1=0$. Điểm nào sau đây thuộc mặt phẳng $(P)$?
	 \choice
		{\True  $A(1;1;3)$}
		{$B(1;1;-3)$} 
		{$C(3;1;1)$} 
		{$D(-1;-1;3)$} 
	\loigiai{ 
		Thay tọa độ các điểm vào phương trình  mặt phẳng $(P)\colon x+y-z+1=0$, ta có
		\begin{itemize}
			\item Với $A(1;1;3)\Rightarrow 1+1-3+1=0\Rightarrow 0=0$,  đúng;
			 \item Với $B(1;1;-3)\Rightarrow 1+1+3+1=6\ne 0$, sai;
			\item Với $C(3;1;1)\Rightarrow  3+1-1+1=4\ne 0$, sai;
			\item Với $D(-1;-1;3)\Rightarrow-1-1-3+1=-4\ne 0$, sai. 
		\end{itemize}
	} 
\end{ex}

% nguồn: Phan1_Lop_12_Hinh_hoc_Chuong_5
\begin{ex}%[2H5N1-2]%[ma:f2507b1c]
	Trong không gian $Oxyz$, véc-tơ nào dưới đây là một véc-tơ pháp  tuyến của mặt phẳng $(P)\colon 2x-y-5z+1=0$?
	 \choice[2]
		{\True $\overrightarrow{n}_1=(2;-1;-5)$} 
		{$\overrightarrow{n}_2=(2;1;-5)$} 
		{$\overrightarrow{n}_3=(2;1;5)$} 
		{$\overrightarrow{n}_4=(2;-1;5)$} 
	\loigiai{Mặt phẳng $(P)\colon 2x-y-5z+1=0$ có véc-tơ pháp tuyến là  $\overrightarrow{n}_1=\left(2;-1;-5\right)$.}
\end{ex}

% nguồn: Phan1_Lop_12_Hinh_hoc_Chuong_5
\begin{ex}%[2H5N1-2]%[ma:a3e1c504]
	Trong  không gian $Oxyz$, cho mặt phẳng $(P)$ có phương trình $x - 2y - 5z + 7 = 0$. Một vectơ pháp tuyến của mặt phẳng $(P)$ là
	\choice 
		{$\overrightarrow{n}=(-1;2;-5)$} 
		 {\True $\overrightarrow{n}=(1;-2;-5)$}
		 {$\overrightarrow{n}=(1;2;5)$}
		 {$\overrightarrow{n}=(-2;5;-1)$}
	\loigiai{Một vectơ pháp tuyến của  mặt phẳng $(P)$ là $\overrightarrow{n}=(1;-2;-5)$.}
\end{ex}

% nguồn: Phan1_Lop_12_Hinh_hoc_Chuong_5
\begin{ex}%[2H5N1-2]%[ma:e0f687b3]
	Trong không gian $Oxyz$,  cho mặt phẳng $(P)\colon 2x+3y+z+2=0$. Vectơ nào dưới đây là một vectơ pháp tuyến của $(P)$?
	 \choice
		{$\vec{n}_4=(2;0;3)$} 
		 {$\vec{n}_3=(2;3;2)$}
		{\True $\vec{n}_2=(2;3;1)$} 
		{$\vec{n}_1=(2;3;0)$} 
	\loigiai{Vectơ pháp  tuyến của $(P) \colon 2x+3y+z+2=0$ là $\vec{n}=(2;3;1)$.}
\end{ex}

% nguồn: Phan1_Lop_12_Hinh_hoc_Chuong_5
\begin{ex}%[2H5N1-2]%[ma:7a5c66df]
	Trong  không gian với hệ tọa độ $Oxyz$, cho mặt phẳng $(P)\colon 2x - y + z + 5 = 0$.
	Một vectơ  pháp tuyến của mặt phẳng $(P)$ là
	 \choice
		 {$\overrightarrow{n} = (2; 1; 1)$}
		{$\overrightarrow{n} = (-2; 1; -1)$} 
		{\True  $\overrightarrow{n} = (2; -1; 1)$}
		{$\overrightarrow{n} =  (-2; -1; 1)$}
	 \loigiai{
		Véctơ pháp tuyến của mặt phẳng $Ax + By +  Cz + D = 0$ là $\overrightarrow{n} = (A,B,C)$.\\
		Vậy với $(P): 2x - y + z + 5 =  0$, vectơ pháp tuyến là $\overrightarrow{n} = (2; -1; 1)$.
	} 
\end{ex}

% nguồn: Phan1_Lop_12_Hinh_hoc_Chuong_5
\begin{ex}%[2H5N1-2]%[ma:aa01ce08]
	Trong không  gian $Oxyz$, mặt phẳng $(P)\colon2x-y-3=0$ có một vectơ pháp tuyến là
	 \choice
		 {$\vec{n}_1=\left(2;-1\right)$}
		 {$\vec{n}_2=\left(2;-1;-3\right)$}
		{$\vec{n}_3=\left(2;-1;0\right)$} 
		{\True  $\vec{n}_4=\left(-2;1;3\right)$}
	\loigiai{ 
		Ta có $\vec{n}_3=\left(2;-1;0\right)$ là $1$ vectơ pháp tuyến của mặt phẳng $(P)$. 
	} 
\end{ex}

% nguồn: Phan1_Lop_12_Hinh_hoc_Chuong_5
\begin{ex}%[2H5N1-2]%[ma:a6602e2d]
	 Trong không gian với hệ tọa độ $Oxyz$, cho mặt phẳng $(P)\colon 3x+2y-z+1=0$. Mặt phẳng $(P)$ có vectơ pháp tuyến là
	 \choice
		{$\overrightarrow{n}=(3;-1;2)$} 
		{\True $\overrightarrow{n}=(-1;3;2)$} 
		{$\overrightarrow{n}=(3;2;-1)$} 
		{$\overrightarrow{n}=(2;3;-1)$} 
	\loigiai{Mặt phẳng $(P)$ có vectơ pháp tuyến  là $\overrightarrow{n}=(3; 2;-1)$.}
\end{ex}

% nguồn: Phan1_Lop_12_Hinh_hoc_Chuong_5
\begin{ex}%[2H5N1-2]%[ma:ea0e6e6a]
	 Trong không gian $Oxyz$, vectơ nào dưới đây là một vectơ pháp tuyến của mặt phẳng $(Oxy)$?
	\choice 
		{$\overrightarrow{n}_1 =  (1;0;1)$}
		{$\overrightarrow{n}_2 =  (1;1;0)$}
		{\True $\overrightarrow{n}_3 = (0;0;1)$} 
		{$\overrightarrow{n}_4 = (1;0;0)$} 
	 \loigiai{
		Mặt phẳng $(Oxy)$ có vec-tơ pháp tuyến  là $\overrightarrow{n} = (0;0;1)$.
	 }
\end{ex}

% nguồn: Phan1_Lop_12_Hinh_hoc_Chuong_5
\begin{ex}%[2H5N1-2]%[ma:072c375a]
	Trong không gian với hệ trục tọa độ $Oxyz$, cho mặt phẳng $(P)$ có phương trình $4x-2y+2025=0$. Véc-tơ  nào sau đây là một véc-tơ pháp tuyến của mặt phẳng $(P)$?
	\choice 
		{$\overrightarrow{n_1}=(4;-2;2025)$} 
		{\True $\overrightarrow{n_3}=(-2;1;0)$} 
		 {$\overrightarrow{n_4}=(1;2;0)$}
		{$\overrightarrow{n_2}=(2;1;0)$} 
	 \loigiai{
		Mặt  phẳng $(P)$ có véc-tơ pháp tuyến là $\overrightarrow{n_3}=-2(-2;1;0)$.
	 }
\end{ex}

% nguồn: Phan1_Lop_12_Hinh_hoc_Chuong_5
\begin{ex}%[2H5N1-2]%[ma:86e90b32]
	Trong không gian Oxyz , cho mặt phẳng $(P)\colon x-2y+3=0$. véc-tơ  nào dưới đây là một véc-tơ
	pháp tuyến của của mặt phẳng  $(P)$?
	\choice 
		{$\overrightarrow{n}_3=(1;-2;3)$} 
		{\True  $\overrightarrow{n}_1=(1;-2;0)$}
		{$\overrightarrow{n}_2=(1;0;-2)$} 
		 {$\overrightarrow{n}_4=(-1;0;2)$}
	 \loigiai{
		Một véc-tơ 
		pháp tuyến  của của mặt phẳng $(P)\colon x-2y+3=0$ là $\overrightarrow{n}_1=(1;-2;0)$.
	} 
\end{ex}

% nguồn: Phan1_Lop_12_Hinh_hoc_Chuong_5
\begin{ex}%[2H5N1-2]%[ma:684044db]
	Trong không gian $Oxyz$, mặt phẳng $(P)\colon x+2y+3z-5=0$  có một vectơ pháp tuyến là
	 \choice
		{$\vec{n}=(1;  2;-3)$}
		{\True $\vec{n}=(1; 2;  3)$}
		{$\vec{n}=(3; 2; 1)$} 
		{$\vec{n}=(-1; 2;  3)$}
	\loigiai{ 
		Mặt phẳng $(P)$ có một vectơ pháp tuyến  là $\vec{n}=(1; 2; 3)$.
	 }
\end{ex}

% nguồn: Phan1_Lop_12_Hinh_hoc_Chuong_5
\begin{ex}%[2H5N1-2]%[ma:5a9075d9]
	Trong không gian $Oxyz$, véc-tơ nào dưới đây là một véc-tơ pháp tuyến  của mặt phẳng $\left(P\right)\colon 2x-y-5z+1=0$?
	\choice 
		{\True $\vec{n}_1=\left(2;-1;-5\right)$} 
		{$\vec{n}_2=\left(2;1;-5\right)$} 
		 {$\vec{n}_3=\left(2;1;5\right)$}
		 {$\vec{n}_4=\left(2;-1;5\right)$}
	\loigiai{Mặt phẳng $\left(P\right)\colon 2x-y-5z+1=0$ có véc-tơ pháp tuyến  là $\vec{n}_1=\left(2;-1;-5\right)$.}
\end{ex}

% nguồn: Phan1_Lop_12_Hinh_hoc_Chuong_5
\begin{ex}%[2H5N1-2]%[ma:2a8c221b]
	Trong  không gian $Oxyz$, cho mặt phẳng $(P) \colon 2x-z+1=0$. Vectơ nào dưới đây~là một vectơ pháp tuyến của $(P)$?
	 \choice
		 {$\vec{n}_1=(2;-1;1)$}
		 {$\vec{n}_2=(2;0;1)$}
		{\True  $\vec{n}_3=(2;0;-1)$}
		{$\vec{n}_4=(2;-1;0)$} 
	 \loigiai{
	 Ta có $\vec{n}_{(P)}=(2;0;-1)=\vec{n}_3$}
\end{ex}

% nguồn: Phan1_Lop_12_Hinh_hoc_Chuong_5
\begin{ex}%[2H5N1-2]%[ma:d0b1bb5a]
	Trong không gian $Oxyz$, mặt phẳng  $(P) \colon x+y+z-1=0$ có một vectơ pháp tuyến là
	 \choice
		{$\vec{n}_{1}  =(-1;1;1)$}
		{$\vec{n}_{2} = (1;-1;1)$} 
		{\True $\vec{n}_{3}  = (1;1;1)$}
		{$\vec{n}_{4} =  (1;1;-1)$}
	\loigiai{ 
		Một vec-tơ pháp tuyến của mặt phẳng $(P):x+y+z-1=0$  là $\vec{n}_{3} = (1;1;1)$.
	} 
\end{ex}

% nguồn: Phan1_Lop_12_Hinh_hoc_Chuong_5
\begin{ex}%[2H5N1-2]%[ma:c99c486d]
	Cho mặt phẳng $(\alpha)$ có phương  trình là $x+2 y-1=0$, vectơ pháp tuyến của mặt phẳng $(\alpha)$ là?
	 \choice
		{$\vec{n}=(1;  2;-1)$}
		{$\vec{n}=(-1;-2; 1)$} 
		{\True  $\vec{n}=(1; 2; 0)$}
		 {$\vec{n}=(1;-2; 0)$}
	 \loigiai{
		Vectơ  pháp tuyến của mặt phẳng $(\alpha)$ được xác định trực tiếp từ phương trình là $\vec{n}=(1;2;0)$.
	 }
\end{ex}

% nguồn: Phan1_Lop_12_Hinh_hoc_Chuong_5
\begin{ex}%[2H5N1-2]%[ma:9c6ab657]
	Trong không gian với hệ tọa độ $Oxyz$, mặt phẳng  $2x - 4y + 8z - 1 = 0$ có một vectơ pháp tuyến là
	 \choice
		{$\overrightarrow{n}=(1; 2;  4)$}
		{\True $\overrightarrow{n}=(1;  -2; 4)$}
		{$\overrightarrow{n}=(2;  -4; 4)$}
		{$\overrightarrow{n}=(2;  -4; 1)$}
	\loigiai{ 
		Mặt phẳng $2x - 4y + 8z - 1 = 0$ có một vectơ pháp tuyến là $\overrightarrow{n}=(1;  -2; 4)$.
	 }
\end{ex}

% nguồn: Phan1_Lop_12_Hinh_hoc_Chuong_5
\begin{ex}%[2H5N1-2]%[ma:d4bbbe47]
	[GV:Xuan  Vy Pham]
	Trong không gian $Oxyz$, cho mặt phẳng $(P)$ vuông góc với đường thẳng \\$d \colon  \dfrac{x-2}{3}=\dfrac{y}{1}=\dfrac{z+2}{-2}$. Vectơ nào dưới đây là một vectơ pháp tuyến của $(P)$?
	\choice 
		{$\overrightarrow{n}=\left(-3;-1;-2\right)$} 
		 {\True $\overrightarrow{n}=\left(3;1;-2\right)$}
		{$\overrightarrow{n}=\left(3;2;1\right)$} 
		 {$\overrightarrow{n}=\left(3;-1;2\right)$}
	\loigiai{Đường thẳng $d \colon \dfrac{x-2}{3}=\dfrac{y}{1}=\dfrac{z+2}{-2}$ có một  vectơ chỉ phương là $\overrightarrow{u_d}=(3;1;-2)$.\\
		 Vì mặt phẳng $(P)$ vuông góc với đường thẳng $d$ nên một vectơ pháp tuyến của $(P)$ là vectơ chỉ phương $\overrightarrow{u_d}$ của đường thẳng $d$.\\
	 Do đó $\overrightarrow{n}=\overrightarrow{d}=\left(3;1;-2\right)$.}
\end{ex}

% nguồn: Phan1_Lop_12_Hinh_hoc_Chuong_5
\begin{ex}%[2H5N1-2]%[ma:7187f10c]
	Trong không gian $Oxyz$, cho mặt phẳng $(P):x-2y+3=0$. Vectơ nào dưới đây là  một vectơ pháp tuyến của $(P)$?
	 \choice
		{$\vec{n}=(1;-2;3)$} 
		{\True $\vec{n}=(-1;2;0)$} 
		{$\vec{n}=(1;0;-2)$} 
		 {$\vec{n}=(3;-1;0)$}
	 \loigiai{
		$(P)$ có một vectơ pháp tuyến $\vec{a}=(1;-2;0)$ nên $\vec{n}=-\vec{a}=(-1;2;0)$ là một  vectơ pháp tuyến của $(P)$.
	} 
\end{ex}

% nguồn: Phan1_Lop_12_Hinh_hoc_Chuong_5
\begin{ex}%[2H5N1-2]%[ma:7947e174]
	Trong không gian  với hệ tọa độ $Oxyz$, mặt phẳng nào dưới đây nhận\break $\overrightarrow{n} = (3;1;-7)$ là một vectơ pháp tuyến?
	\choice 
		{$3x + z  + 7 = 0$}
		{$3x - y - 7z + 1 = 0$} 
		{$3x + y - 7 =  0$}
		{\True $3x + y - 7z - 3 =  0$}
	 \loigiai{
		Phương trình  mặt phẳng $3x + y - 7z - 3 = 0$ có vectơ pháp tuyến là $\overrightarrow{n} = (3;1;-7)$.
	 }
\end{ex}

% nguồn: Phan1_Lop_12_Hinh_hoc_Chuong_5
\begin{ex}%[2H5N1-2]%[ma:9bf0e3f4]
	Trong không gian  $Oxyz$, cho mặt phẳng $(P)$ có phương trình $-2 x+2 y-z-2025=0$. Mặt phẳng $(P)$ có vectơ pháp tuyến là
	\choice 
		{\True $(-2;  2;-1)$}
		{$(-2; 2;-3)$} 
		 {$(-4; 4; 2)$}
		 {$(0; 0;-3)$}
	\loigiai{ 
		 Dựa vào phương trình mặt phẳng $(P)$, ta có vectơ pháp tuyến là $(-2; 2;-1)$.
	 }
\end{ex}

% nguồn: Phan1_Lop_12_Hinh_hoc_Chuong_5
\begin{ex}%[2H5N1-2]%[ma:b3cbe657]
	Trong không gian  với hệ tọa độ $Oxyz$, cho mặt phẳng $(P)\colon 3x+2y-z+1=0$. Mặt phẳng $(P)$ có véc-tơ pháp tuyến là
	\choice 
		 {$\overrightarrow{n}=(3;-1;2)$}
		{$\overrightarrow{n}=(-1;3;2)$} 
		{\True$\overrightarrow{n}=(3;2;-1)$} 
		 {$\overrightarrow{n}=(2;3;-1)$}
	\loigiai{Mặt phẳng $(P)$ có véc-tơ pháp  tuyến là $\overrightarrow{n}=(3; 2;-1)$.}
\end{ex}

% nguồn: Phan1_Lop_12_Hinh_hoc_Chuong_5
\begin{ex}%[2H5N1-2]%[ma:b07c0e7b]
	\immini{Cho hình hộp $ABCD.A'B'C'D'$. Cặp vectơ nào sau đây là cặp vectơ chỉ phương của mặt phẳng $(ABCD)$? 
	 \choice
		 {$\overrightarrow{AB}$, $\overrightarrow{AA'}$}
		{$\overrightarrow{AB}$, $\overrightarrow{AB'}$} 
		{$\overrightarrow{AB}$,  $\overrightarrow{A'B'}$}
		{\True $\overrightarrow{AD}$,  $\overrightarrow{A'B'}$}}
		{\resizebox{0.35\linewidth}{!}{\begin{tikzpicture}
		 \coordinate (x) at (-2,-2);
		\coordinate (y) at (6,0); 
		\coordinate (z)  at (1,3);
		\coordinate (A) at (0,0); 
		\coordinate  (B) at ($(A)+(x)$);
		\coordinate  (C) at ($(B)+(y)$);
		 \coordinate (D) at ($(C)-(x)$);
		\coordinate (A')  at ($(A)+(z)$);
		\coordinate (B')  at ($(B)+(z)$);
		\coordinate  (C') at ($(C)+(z)$);
		\coordinate (D')  at ($(D)+(z)$);
		\draw[dashed]  (B)--(A)--(D)
			(B')--(A)--(A'); 
		\draw (D)--(C) 
			 (D)--(D');
		\draw  (B')--(C')--(D')--(A')--cycle;
		 \draw (B')--(B)--(C)--(C');
		 \foreach \t/\g in {A/-90,B/180,C/0,D/0,A'/90,B'/145,C'/90,D'/90}{
			\fill (\t)  circle (1pt) node[shift={(\g:9pt)}]{$\t$};
		 }
	\end{tikzpicture}
		 }}
	\loigiai{Cặp vectơ chỉ phương của mặt phẳng $(ABCD)$ là $2$ vectơ không cùng phương, có giá song song hoặc nằm trên mặt phẳng $(ABCD)$ đó là $\overrightarrow{AD}$ và $\overrightarrow{A'B'}$. Vì $\overrightarrow{AD}$ có giá nằm trong mặt phẳng $(ABCD)$ và $\overrightarrow{A'B'}$ có giá song  song với mặt phẳng $(ABCD)$.}
\end{ex}

% nguồn: Phan1_Lop_12_Hinh_hoc_Chuong_5
\begin{ex}%[2H5N1-2]%[ma:ac5659c1]
	Trong không gian $Oxyz$, cho mặt phẳng $(P) \colon x+2y-3z+4=0$. Vectơ nào sau đây là một vectơ pháp tuyến của  mặt phẳng $(P)$?
	\choice 
		 {\True $\overrightarrow{n}_1=(1;2;-3)$}
		 {$\overrightarrow{n}_2=(1;2;3)$}
		 {$\overrightarrow{n}_3=(1;-2;3)$}
		{$\overrightarrow{n}_4=(-1;2;3)$} 
	 \loigiai{
		Ta có  vectơ pháp tuyến của $(P)$ là $\overrightarrow{n}_1=(1;2;-3)$.
	 }
\end{ex}

% nguồn: Phan1_Lop_12_Hinh_hoc_Chuong_5
\begin{ex}%[2H5N1-2]%[ma:021833aa]
	Trong không  gian $Oxyz$, mặt phẳng $(P)$ đi qua điểm $A(1;-1;2)$ và nhận $\overrightarrow{n}=(2;1;2)$ làm vectơ pháp tuyến có phương trình là
	\choice 
		{$2x+y+2z+3=0$} 
		 {$x-y+2z-5=0$}
		 {$x-y+2z+3=0$}
		 {\True $2x+y+2z-5=0$}
	 \loigiai{
		 Phương trình mặt phẳng $(P)$ là $2(x-1)+1(y+1)+2(z-2)=0\Leftrightarrow 2x+y+2z-5=0$.
	 }
\end{ex}

% nguồn: Phan1_Lop_12_Hinh_hoc_Chuong_5
\begin{ex}%[2H5N1-2]%[ma:0b8506fe]
	Trong không  gian với hệ tọa độ $Oxyz$, cho mặt phẳng $(P)\colon 2x + 3y + z + 2 = 0$. Vectơ nào dưới đây là một vectơ pháp tuyến của $(P)$?
	 \choice
		{\True  $\overrightarrow{n}_1=(2;3;1)$}
		{$\overrightarrow{n}_2=(2;3;2)$} 
		 {$\overrightarrow{n}_3=(2;3;0)$}
		{$\overrightarrow{n}_4=(2;0;3)$} 
	\loigiai{ 
		Vectơ pháp tuyến của  $(P)$ là $\overrightarrow{n}_P=(2;3;1)$.
	 }
\end{ex}

% nguồn: Phan1_Lop_12_Hinh_hoc_Chuong_5
\begin{ex}%[2H5N1-2]%[ma:50d55816]
	[GV:Xuan Vy Pham] 
	Mặt phẳng $(\alpha) \colon 2x-5y-z+1=0$ có một vectơ pháp  tuyến là
	\choice 
		{$\overrightarrow{n}=(2;5;1)$} 
		 {$\overrightarrow{n}=(2;5;-1)$}
		{\True $\overrightarrow{n}=(-4;10;2)$} 
		 {$\overrightarrow{n}=(-2;5;-1)$}
	\loigiai{Mặt phẳng $(\alpha) \colon 2x-5y-z+1=0$ có  một vectơ pháp tuyến là $(2;-5;-1)$.\\
	Vì $\overrightarrow{n}=(-4;10;2)=-2(2;-5;-1)$  nên $\overrightarrow{n}=(-4;10;2)$ là một vectơ pháp tuyến của mặt phẳng $(\alpha)$.}
\end{ex}

% nguồn: Phan1_Lop_12_Hinh_hoc_Chuong_5
\begin{ex}%[2H5N1-2]%[ma:49c6583e]
	Trong không gian $Oxyz$, cho mặt phẳng $(P)$  có phương trình $-2x+2y-z-3=0$. Mặt phẳng $(P)$ có vectơ pháp tuyến là
	 \choice
		{\True $(0;  0;-3)$}
		 {\True $(4;-4; 2)$}
		 {$(-2; 2;-3)$}
		{$(-4;  4; 2)$}
	\loigiai{ 
		Dựa vào phương trình mặt phẳng $(P)$, ta được vectơ pháp tuyến của $(P)$ là $(4; -4;2)$ vì nó cùng phương với $(-2;  2;-1)$.
	} 
\end{ex}

% nguồn: Phan1_Lop_12_Hinh_hoc_Chuong_5
\begin{ex}%[2H5N1-2]%[ma:0d9abb3c]
	Trong không gian $Oxyz$, cho mặt phẳng $(P)\colon 2x+3y+z+1=0$. Một VTPT của mặt phẳng $(P)$ là 
	\choice 
		 {\True $(2;3;1)$}
		{$(-2;-3;1)$} 
		 {$(2;-3;1)$}
		 {$(3;-2;1)$}
	 \loigiai{
		Vì $(P)\colon 2x+3y+z+1=0$ nên một  VTPT của mặt phẳng $(P)$ là $(2;3;1)$.
	} 
\end{ex}

% nguồn: Phan1_Lop_12_Hinh_hoc_Chuong_5
\begin{ex}%[2H5N1-2]%[ma:c4b2a37b]
	Trong không gian $Oxyz$, cho mặt phẳng $(P)\colon x-3y-4z+1=0$. Một vectơ pháp tuyến của mặt phẳng $(P)$ có  tọa độ là
	\choice 
		{$(1;-3;4)$} 
		{$(1;3;4)$} 
		 {$(-1;-3;4)$}
		{\True $(1;-3;-4)$} 
	\loigiai{ 
		Ta có mặt phẳng $(P)\colon  x-3y-4z+1=0$ nên mặt $(P)$ có một vectơ pháp tuyến là $(1;-3;-4)$.
	 }
\end{ex}

% nguồn: Phan1_Lop_12_Hinh_hoc_Chuong_5
\begin{ex}%[2H5N1-2]%[ma:85448648]
	Trong không gian với hệ tọa độ $Oxyz$, vectơ nào sau đây là một vectơ pháp tuyến của mặt  $y-2x-5z+8=0$
	\choice 
		{$\overrightarrow{n}_1=(1;2;5)$} 
		 {$\overrightarrow{n}_2=(1;-2;-5)$}
		 {\True $\overrightarrow{n}_3=(-2;1;-5)$}
		 {$\overrightarrow{n}_4=(2;1;5)$}
	\loigiai{Ta có $y-2x-5z+8=0 \Leftrightarrow -2x+y-5z+8=0$ suy ra mặt phẳng có một vectơ  pháp tuyến là $\overrightarrow{n}_3=(-2;1;-5)$.}
\end{ex}

% nguồn: Phan1_Lop_12_Hinh_hoc_Chuong_5
\begin{ex}%[2H5N1-2]%[ma:3a1a27e7]
	Trong không gian $Oxyz$, cho mặt phẳng $(P)$ có phương trình $2x+y-z+9=0$. Véc-tơ nào sau đây là  một véc-tơ pháp tuyến của mặt phẳng $(P)$?
	 \choice
		{\True $\overrightarrow{n}_{1}=(2;1;-1)$} 
		{$\overrightarrow{n}_{2}=(2;1;1)$} 
		{$\overrightarrow{n}_{3}=(2;0;0)$} 
		 {$\overrightarrow{n}_{4}=(2;-1;6)$}
	 \loigiai{
		 Một véc-tơ pháp tuyến của mặt phẳng $(P)\colon2x+y-z+9=0$ là $\overrightarrow{n}_{1}=(2;1;-1)$.
	} 
\end{ex}

% nguồn: Phan1_Lop_12_Hinh_hoc_Chuong_5
\begin{ex}%[2H5H1-2]%[ma:2c1d4dd6]
	Trong không gian với hệ tọa độ $Oxyz$, cho điểm $A(1;2;3)$. Mặt phẳng chứa điểm $A$ và  trục $Oz$ có phương trình là
	\choice 
		 {\True $2x - y = 0$}
		{$x + y - z =  0$}
		{$3y - 2z = 0$} 
		{$3x - z = 0$} 
	\loigiai{ 
		Gọi $(P)$  là mặt phẳng chứa điểm $A$ và trục $Oz$.\\
		$Oz$ đi qua điểm $O(0;0;0)$ và có vectơ chỉ phương $\overrightarrow{u}  = (0;0;1)$.\\
		Mặt phẳng $(P)$ có vectơ pháp tuyến $\overrightarrow{n} =  [\overrightarrow{u}, \overrightarrow{OA}] = (2;-1;0)$ và đi qua điểm $O(0;0;0)$ nên có phương trình là
		$$2x - y =  0.$$
	} 
\end{ex}

% nguồn: Phan1_Lop_12_Hinh_hoc_Chuong_5
\begin{ex}%[2H5H1-2]%[ma:23519b4e]
	Trong không gian $Oxyz$, cho ba điểm $A(2;-1;3), B(4;0;1), C(-10;5;3)$. Một cặp vectơ chỉ  phương của mặt phẳng $(ABC)$~là
	 \choice
		{$\overrightarrow{AB}=(2;1;-2),  \overrightarrow{AC}=(-12;6;0)$}
		{$\overrightarrow{AB}=(2;1;2),  \overrightarrow{AC}=(12;6;0)$}
		 {\True $\overrightarrow{AB}=(-2;1;2), \overrightarrow{AC}=(-12;6;0)$}
		{$\overrightarrow{AB}=(2;1;-2), \overrightarrow{AC}=(-12;6;3)$} 
	 \loigiai{
		Ta có  $\heva{&\overrightarrow{AB}=(-2;1;-2)\\&\overrightarrow{AC}=(-12;6;0).}$
	} 
\end{ex}

% nguồn: Phan1_Lop_12_Hinh_hoc_Chuong_5
\begin{ex}%[2H5H1-2]%[ma:b185374a]
	Trong không  gian $Oxyz$, cho ba điểm $A(2;1;-1), B(3;0;1), C(2;-1;3)$. Một vectơ pháp tuyến của mặt phẳng $(ABC)$~là
	 \choice
		 {\True $\vec{n_1}=(0;-4;-2)$}
		 {$\vec{n_2}=(0;4;-2)$}
		{$\vec{n_3}=(1;-1;2)$} 
		{$\vec{n_4}=(0;-2;4)$} 
	 \loigiai{
		Mặt phẳng $(ABC)$ có vectơ pháp  tuyến $\vec{n}$ vuông góc với $\overrightarrow{AB}=(1;-1;2), \overrightarrow{BC}=(-1;-1;2)$
		$$\Rightarrow \vec{n}=\left[\overrightarrow{AB},  \overrightarrow{BC}\right]=(0;-4;-2)$$
	} 
\end{ex}

% nguồn: Phan1_Lop_12_Hinh_hoc_Chuong_5
\begin{ex}%[2H5H1-2]%[ma:fd7a59c6]
	Cho mặt phẳng $(P)\colon x+2 y+3 z-1=0$. Vectơ nào dưới đây là một vectơ  pháp tuyến của $(P)$?
	 \choice
		{$\overrightarrow{n}_1=(1;  3;-1)$}
		{$\overrightarrow{n}_2=(2; 3;-1)$} 
		{$\overrightarrow{n}_3=(1;  2;-1)$}
		{\True  $\overrightarrow{n}_4=(1; 2; 3)$}
	 \loigiai{
		$(P)$ có  vectơ pháp tuyến là $\overrightarrow{n}_4=(1; 2; 3)$.
	 }
\end{ex}

% nguồn: Phan1_Lop_12_Hinh_hoc_Chuong_5
\begin{ex}%[2H5V1-2]%[ma:22d1372e]
	Trong không gian $Oxyz$, cho hai điểm  $A(2;4;1), B(-1;1;3)$ và mặt phẳng \break $(P) \colon x-3y+2z-5=0$. Mặt phẳng $(Q)$ đi qua hai điểm $A, B$ và vuông góc với $(P)$ có một vectơ pháp tuyến~là $\vec{n}=(a;b;c)$ với $c$~là số nguyên tố. Giá trị của $a+b+c$ bằng
	 \choice
		{$9$} 
		{$3$} 
		{\True $5$} 
		 {$7$}
	\loigiai{ 
		Ta  có $\vec{n}_{(Q)}\bot\overrightarrow{AB}$ và $\vec{n}_{(Q)}\bot\vec{n}_{(P)}$
		$$\Rightarrow  \vec{n}_{(Q)}=\left[\overrightarrow{AB},\vec{n}_{(P)}\right]=(0;8;12)$$
		Ta  có $\vec{n}=(0;2;3)$~là vectơ cùng phương với $\vec{n}_{(Q)}$ thỏa mãn yêu cầu.\\
		 Suy ra $a+b+c=5$
	.} 
\end{ex}

% nguồn: Phan1_Lop_12_Hinh_hoc_Chuong_5
\begin{ex}%[2H5B1-2]%[ma:c92f50a9]
	Trong không gian với hệ tọa~độ $Oxyz$, mặt phẳng nhận $\vec{n}=\left(a;b;c \right)$~là vectơ pháp tuyến và đi  qua điểm $I\left({{x}_0};{{y}_0};{{z}_0} \right)$ có phương trình~là
	 \choice
		{$a\left(x-{{x}_0} \right)+b\left(y-{{y}_0}  \right)+c\left(z+{{z}_0} \right)=0$}
		{\True $a\left(x-{{x}_0} \right)+b\left(y-{{y}_0} \right)+c\left(z-{{z}_0}  \right)=0$}
		{$a\left(x-{{x}_0} \right)+b\left(y-{{y}_0}  \right)-c\left(z-{{z}_0} \right)=0$}
		 {$a\left(x-{{x}_0} \right)+c\left(y-{{y}_0} \right)+b\left(z-{{z}_0} \right)=0$}
	\loigiai{ 
		Mặt phẳng  nhận $\vec{n}=\left(a;b;c \right)$~là vectơ pháp tuyến và đi qua điểm $I\left({{x}_0};{{y}_0};{{z}_0} \right)$ có phương trình:$a\left(x-{{x}_0} \right)+b\left(y-{{y}_0} \right)+c\left(z-{{z}_0} \right)=0$.
	 }
\end{ex}

% nguồn: Phan1_Lop_12_Hinh_hoc_Chuong_5
\begin{ex}%[2H5B1-2]%[ma:db033285]
	Trong không gian với hệ tọa~độ $Oxyz$, cho hình hộp $ABCD.A'B'C'D'$. Cặp vectơ nào~là cặp vectơ chỉ  phương của mặt phẳng $\left(ABCD \right)$?
	 \choice
		{\True$\overrightarrow{AB}$,  $\overrightarrow{A'D'}$}
		{$\overrightarrow{AB}$,  $\overrightarrow{AD'}$}
		{$\overrightarrow{AB}$, $\overrightarrow{AA'}$} 
		 {$\overrightarrow{AB}$, $\overrightarrow{AB'}$}
	\loigiai{ 
		Trong hình hộp chữ nhật $ABCD.A'B'C'D'$ thì $\overrightarrow{AB}$ và  $\overrightarrow{A'D'}$~là cặp vectơ chỉ phương của $(ABCD)$.
	} 
\end{ex}

% nguồn: Phan1_Lop_12_Hinh_hoc_Chuong_5
\begin{ex}%[2H5B1-2]%[ma:886abfd3]
	Trong không gian với hệ tọa~độ $Oxyz$, khoảng cách  từ điểm $M\left({{x}_0};{{y}_0};{{z}_0} \right)$ đến mặt phẳng $(P)\colon ax+by+cz+d=0$~là
	\choice 
		 {$\dfrac{\left| a{{x}_0}+b{{y}_0}+c{{z}_0}+d \right|}{\sqrt{a^2+b^2+c^2+d^2}}$}
		{$\dfrac{\left|  a{{x}_0}+b{{y}_0}-c{{z}_0}+d \right|}{\sqrt{a^2+b^2+c^2}}$}
		{\True $\dfrac{\left| a{{x}_0}+b{{y}_0}+c{{z}_0}+d \right|}{\sqrt{a^2+b^2+c^2}}$} 
		 {$\dfrac{\left| a{{x}_0}+c{{y}_0}+b{{z}_0}+d \right|}{\sqrt{a^2+b^2+c^2}}$}
	\loigiai{ 
		khoảng cách từ điểm $M\left({{x}_0};{{y}_0};{{z}_0} \right)$ đến mặt phẳng $\left(P \right)\colon  ax+by+cz+d=0$~là $$d\left(M,(P)\right)=\dfrac{\left| a{{x}_0}+b{{y}_0}+c{{z}_0}+d \right|}{\sqrt{a^2+b^2+c^2}}.$$
	} 
\end{ex}

% nguồn: Phan1_Lop_12_Hinh_hoc_Chuong_5
\begin{ex}%[2H5B1-2]%[ma:f259fbac]
	Trong không gian với hệ tọa~độ $Oxyz$, cho hình hộp chữ nhật  $ABCD.A'B'C'D'$. Vectơ nào~là vectơ pháp tuyến của mặt phẳng $\left(ABCD \right)$?
	 \choice
		{$\overrightarrow{AB}$} 
		{$\overrightarrow{AC}$} 
		{$\overrightarrow{AD}$} 
		{\True  $\overrightarrow{AA'}$}
	\loigiai{ 
		Trong hình hộp chữ nhật $ABCD.A'B'C'D'$thì $AA'\perp (ABCD)$.  Do đó $\overrightarrow{AA'}$~là một vectơ pháp tuyến của $(ABCD)$.
	 }
\end{ex}


%%%%%%%%%%%%%%%% ĐÚNG/SAI %%%%%%%%%%%%%%%%

% nguồn: Phan1_Lop_12_Hinh_hoc_Chuong_5
\begin{ex}%[2H5H1-2]%[ma:d8077aeb]
	Trong  không gian $Oxyz$, cho các điểm $A(1; -2; -1)$, $B(4; 1; 2)$, $C(2; 3; 1)$.
	 \choiceTF
		{\True $\overrightarrow{AB}  = (3; 3; 3)$}
		{\True Ba điểm  $A, B, C$ không thẳng hàng}
		{\True Mặt phẳng đi qua ba điểm $A,  B, C$ có vectơ pháp tuyến là $\overrightarrow{a} = (3; 1; -4)$}
		{Mặt phẳng $(\alpha)$  đi qua $A$ đồng thời song song với $Oy$ và đường thẳng $BC$ có vectơ pháp tuyến là $\overrightarrow{n} = (1; 0; 2)$}
	 \loigiai{
		\begin{itemchoice}
			 \itemch $\overrightarrow{AB} = (4 - 1; 1 - (-2); 2 - (-1)) = (3; 3; 3)$.
			\itemch $\overrightarrow{AC} = (2 - 1; 3 - (-2); 1 - (-1)) = (1;  5; 2)$.\\
			Vì $\overrightarrow{AB}$ và $\overrightarrow{AC}$ không  cùng phương suy ra ba điểm không thẳng hàng.
			 \itemch $\overrightarrow{a} = [\overrightarrow{AB}, \overrightarrow{AC}] = [(3; 3; 3), (1; 5; 2)] = (-9; -3; 12)$.\\
			Vì $\overrightarrow{a} = (3; 1; -4)$ cùng phương  suy ra cũng là pháp tuyến.
			\itemch $\overrightarrow{j} = (0; 1; 0)$, $\overrightarrow{BC} = (-2; 2; -1)$.\\ 
			Vậy $\overrightarrow{n} = [\overrightarrow{j}, \overrightarrow{BC}] = (-1; 0; 2)$.\\ 
			Suy ra $\overrightarrow{n} = (1; 0;  2)$ không cùng phương.
		\end{itemchoice}
	} 
\end{ex}

% nguồn: Phan1_Lop_12_Hinh_hoc_Chuong_5
\begin{ex}%[2H5H1-2]%[ma:68344600]
	Trong không gian $Oxyz$, cho ba điểm $A(-1; 1; 0)$, $B(1; -1; 2)$, $C(1; -2; 1)$. 
	\choiceTF 
		{\True Một vectơ pháp tuyến của mặt phẳng  $(ABC)$ là $[\overrightarrow{AB}, \overrightarrow{AC}]$}
		 {Vectơ $\overrightarrow{n} = (1; 2; 3)$ là một vectơ pháp tuyến của mặt phẳng $(ABC)$}
		{\True Vectơ $\overrightarrow{u} = (1; 1; 0)$ là một vectơ pháp tuyến  của mặt phẳng đi qua $O$ và chứa đường thẳng $AB$}
		{Vectơ $\overrightarrow{v} = (1; 2; 3)$ là một vectơ pháp tuyến của mặt phẳng  song song với hai đường thẳng $AB$ và $OC$}
	\loigiai{ 
		\begin{itemchoice}
			\itemch Ta có $\overrightarrow{AB} = (2; -2; 2)$, $\overrightarrow{AC} = (2;  -3; 1)$ suy ra $[\overrightarrow{AB}, \overrightarrow{AC}] = (4; 2; -2)$ là một vectơ pháp tuyến của mặt phẳng $(ABC)$.
			\itemch  Xét $\overrightarrow{n} = (1; 2; 3)$ không cùng phương với $[\overrightarrow{AB}, \overrightarrow{AC}] = (4; 2; -2)$ vì $\dfrac{4}{1} \neq \dfrac{2}{2}$.
			\itemch Ta có $\overrightarrow{OA} = (-1; 1; 0)$, $\overrightarrow{OB} = (1; -1; 2)$ suy ra $[\overrightarrow{OA}, \overrightarrow{OB}]  = (2; 2; 0)$.\\
			Vectơ $\overrightarrow{u} = (1; 1;  0)$ cùng phương với $[\overrightarrow{OA}, \overrightarrow{OB}]$.
			\itemch Ta có $\overrightarrow{AB} = (2; -2; 2)$, $\overrightarrow{OC} = (1; -2; 1)$ suy ra $[\overrightarrow{AB}, \overrightarrow{OC}] = (2; 0; -2)$.\\ 
			Vì $\overrightarrow{v} = (1; 2;  3)$ không cùng phương với $[\overrightarrow{AB}, \overrightarrow{OC}]$.
		\end{itemchoice}
	} 
\end{ex}

% nguồn: Phan1_Lop_12_Hinh_hoc_Chuong_5
\begin{ex}%[2H5H1-2]%[ma:0cb37c70]
	Trong  không gian $Oxyz$, cho tam giác $ABC$ với $A(1; 1; 1)$, $B(1; 2; 2)$, $C(4; 1; 0)$.
	 \choiceTF
		{\True $\overrightarrow{AB} =  (0; 1; 1)$}
		{\True Tích có hướng của hai vectơ $\overrightarrow{AB}, \overrightarrow{AC}$ là $\overrightarrow{a}  = (-1; 3; -3)$}
		{$\overrightarrow{BC}, \overrightarrow{b} = (6; -2; -4)$ là cặp  vectơ chỉ phương của mặt phẳng $(ABC)$}
		{Vectơ pháp  tuyến của mặt phẳng $(AOB)$ là $\overrightarrow{n} = (1; 1; 2)$}
	\loigiai{ 
		\begin{itemchoice}
			\itemch Ta có $\overrightarrow{AB} = (1 - 1; 2 - 1; 2 - 1)  = (0; 1; 1)$.
			\itemch Ta có $\overrightarrow{AB} =  (0; 1; 1)$, $\overrightarrow{AC} = (3; 0; -1)$ suy ra
			 $\overrightarrow{a} = [\overrightarrow{AB}, \overrightarrow{AC}] = (-1; 3; -3)$.
			\itemch Ta có $\overrightarrow{BC} = (3; -1; -2)$ và $\overrightarrow{b} = (6; -2; -4) =  2\overrightarrow{BC}$ suy ra $\overrightarrow{BC}, \overrightarrow{b}$ cùng phương nên không phải là cặp vectơ chỉ phương độc lập của mặt phẳng.
			\itemch Ta có $\overrightarrow{OA} = (1; 1; 1)$, $\overrightarrow{OB} = (1; 2; 2)$ suy ra  $\overrightarrow{n} = [\overrightarrow{OA}, \overrightarrow{OB}] = (0; -1; 1)$.\\
			Suy ra $\overrightarrow{n}  = (1; 1; 2)$ không cùng phương.
		\end{itemchoice}
	 }
\end{ex}
